\documentclass[11pt,a4paper]{article}

\usepackage[T1]{fontenc}
\usepackage[utf8]{inputenc}
\usepackage{amsmath,amssymb,amsthm}
\usepackage{mathtools}
\usepackage[margin=2.3cm]{geometry}
\usepackage[protrusion=true,expansion=false]{microtype}
\usepackage{enumitem}
\usepackage[colorlinks=true,linkcolor=black,urlcolor=blue]{hyperref}

\renewcommand{\abstractname}{Résumé}
\renewcommand{\contentsname}{Table des matières}
\renewcommand{\proofname}{Démonstration}

\usepackage{titlesec}
\titleformat{\section}{\normalfont\large\bfseries}{\thesection}{0.7em}{}
\titleformat{\subsection}{\normalfont\bfseries}{\thesubsection}{0.6em}{}
\setlength{\parskip}{0.35em}
\setlength{\parindent}{0pt}

\theoremstyle{plain}
\newtheorem{theoreme}{Théorème}[section]
\newtheorem{lemme}[theoreme]{Lemme}
\newtheorem{proposition}[theoreme]{Proposition}
\newtheorem{corollaire}[theoreme]{Corollaire}
\theoremstyle{definition}
\newtheorem{definition}[theoreme]{Définition}
\theoremstyle{remark}
\newtheorem{remarque}[theoreme]{Remarque}

\newcommand{\R}{\mathbb{R}}
\newcommand{\N}{\mathbb{N}}
\newcommand{\Q}{\mathbb{Q}}
\newcommand{\E}{\mathbb{E}}
\newcommand{\Prob}{\mathbb{P}}
\newcommand{\Var}{\operatorname{Var}}
\newcommand{\Expo}{\mathcal{E}}
\newcommand{\barF}{\bar{F}}
\newcommand{\eqd}{\overset{d}{=}}
\newcommand{\cvloi}{\xrightarrow{\ d\ }}
\newcommand{\Xord}[1]{X_{(#1)}}
\newcommand{\dd}{\,\mathrm{d}}
\newcommand{\Jac}{\mathrm{J}}
\DeclareMathOperator*{\argmin}{arg\,min}

\title{\bfseries Absence de mémoire, taux de panne\\
et statistiques d'ordre exponentielles}
\author{}
\date{\today}

\begin{document}
\maketitle
\vspace{-2.4em}
\begin{center}\rule{0.5\textwidth}{0.4pt}\end{center}

\begin{abstract}
\noindent
On établit d'abord que l'exponentielle est l'unique loi de durée à taux de panne
constant, et que les trois énoncés « absence de mémoire », « survie
multiplicative » et « taux constant » sont équivalents (l'équation de Cauchy
étant résolue en détail). On en déduit ensuite la structure des statistiques
d'ordre d'un échantillon i.i.d.\ exponentiel. Le théorème de Rényi est démontré
deux fois : par changement de variables, avec preuve que l'application des
espacements est un $C^\infty$-difféomorphisme et calcul explicite du jacobien ;
puis par un argument probabiliste de compétition et de renouvellement, rendu
rigoureux par un calcul de loi jointe qui évite tout conditionnement par un temps
aléatoire. Suivent les moments et la convergence vers Gumbel, démontrée
directement.
\end{abstract}

\tableofcontents
\vspace{0.8em}

\begin{center}\rule{0.35\textwidth}{0.4pt}\end{center}
\part*{Partie I --- Survie, taux de panne, absence de mémoire}
\addcontentsline{toc}{section}{\textbf{Partie I --- Survie, taux de panne, absence de mémoire}}

%===================================================================
\section{Fonction de survie}
%===================================================================

\begin{definition}
Soit $X$ une variable aléatoire réelle positive p.s., de fonction de répartition
$F(x)=\Prob(X\le x)$. Sa \emph{fonction de survie} (ou fiabilité) est
\[
\barF(x)\;=\;\Prob(X>x)\;=\;1-F(x),\qquad x\ge 0 .
\]
\end{definition}

$\barF$ est décroissante, continue à droite, $\barF(x)\to 0$ quand
$x\to+\infty$, et $\barF(0)=1$ dès que $\Prob(X=0)=0$. Si $X$ admet une densité
$f$, alors $F$ est absolument continue et $\barF'=-f$ presque partout.

\begin{proposition}[Espérance par la survie]
Pour $X\ge 0$ p.s., $\ \E[X]=\displaystyle\int_0^{\infty}\barF(x)\dd x$, les deux
membres étant simultanément finis ou infinis.
\end{proposition}

\begin{proof}
La fonction $(x,\omega)\mapsto \mathbf 1_{\{x<X(\omega)\}}$ est mesurable positive
sur $(0,\infty)\times\Omega$. Le théorème de Tonelli autorise l'échange :
\[
\int_0^\infty \barF(x)\dd x
=\int_0^\infty \E\bigl[\mathbf 1_{\{X>x\}}\bigr]\dd x
=\E\Bigl[\int_0^\infty \mathbf 1_{\{x<X\}}\dd x\Bigr]
=\E[X].
\]
Aucune hypothèse d'intégrabilité n'est requise, la positivité suffit.
\end{proof}

%===================================================================
\section{Taux de panne et hasard cumulé}
%===================================================================

Dans toute cette section, $X\ge0$ admet une densité $f$ ; on note
$\tau=\sup\{x\ge0:\barF(x)>0\}\in(0,+\infty]$ le bout du support.

\begin{definition}
Pour $x\in[0,\tau)$, le \emph{taux de panne} (fonction de hasard) de $X$ est
\[
h(x)\;=\;\lim_{\varepsilon\downarrow 0}\frac1\varepsilon\,
\Prob\bigl(x<X\le x+\varepsilon \,\bigm|\, X>x\bigr),
\]
lorsque cette limite existe.
\end{definition}

\begin{proposition}[Forme analytique]
Si $f$ est continue en $x\in[0,\tau)$, la limite ci-dessus existe et vaut
\[
h(x)=\frac{f(x)}{\barF(x)}=-\frac{\barF'(x)}{\barF(x)}=-\frac{\dd}{\dd x}\log\barF(x).
\]
\end{proposition}

\begin{proof}
Comme $\barF(x)>0$, le conditionnement est licite et
\[
\frac1\varepsilon\Prob\bigl(x<X\le x+\varepsilon\mid X>x\bigr)
=\frac{1}{\barF(x)}\cdot\frac{F(x+\varepsilon)-F(x)}{\varepsilon}
=\frac{1}{\barF(x)}\cdot\frac1\varepsilon\int_x^{x+\varepsilon} f(u)\dd u .
\]
La continuité de $f$ en $x$ donne $\frac1\varepsilon\int_x^{x+\varepsilon}f\to f(x)$
quand $\varepsilon\downarrow0$. Les deux dernières égalités de l'énoncé résultent
de $\barF'=-f$ et de la dérivation de $\log$.
\end{proof}

\begin{remarque}
$h$ n'est pas une densité : $f(x)$ est un risque rapporté à la population
initiale, $h(x)$ un risque rapporté aux \emph{survivants} à l'instant $x$. En
particulier $h$ n'est pas d'intégrale $1$, et n'est même pas nécessairement
intégrable sur $[0,\tau)$.
\end{remarque}

\begin{theoreme}[Le taux détermine la loi]\label{th:reconstruction}
Soit $X\ge0$ de survie $\barF$ absolument continue sur tout compact de
$[0,\tau)$, avec $\barF(0)=1$. Posons le \emph{hasard cumulé}
$H(x)=\int_0^x h(u)\dd u$. Alors
\[
\barF(x)=\exp\bigl(-H(x)\bigr),\qquad x\in[0,\tau),
\qquad\text{et}\qquad f(x)=h(x)\,e^{-H(x)} .
\]
Réciproquement, toute fonction $h\ge0$ localement intégrable sur $[0,\infty)$
telle que $\int_0^\infty h=+\infty$ définit par cette formule la survie d'une
unique loi sur $[0,\infty)$.
\end{theoreme}

\begin{proof}
Sur $[0,\tau)$ on a $\barF>0$, donc $x\mapsto\log\barF(x)$ est bien définie et
absolument continue, de dérivée $\barF'/\barF=-h$ p.p. Le théorème fondamental
de l'analyse pour les fonctions absolument continues donne
\[
\log\barF(x)-\log\barF(0)=\int_0^x\bigl(\log\barF\bigr)'(u)\dd u=-\int_0^x h(u)\dd u=-H(x),
\]
et $\log\barF(0)=0$, d'où $\barF=e^{-H}$. Puis $f=-\barF'=h\,e^{-H}$.

Réciproquement, soit $h\ge0$ localement intégrable d'intégrale divergente. La
fonction $G=e^{-H}$ est continue, décroissante (car $H$ est croissante),
$G(0)=1$ et $G(x)\to0$ puisque $H(x)\to+\infty$. C'est donc bien une fonction de
survie, et elle détermine $F=1-G$ de façon unique.
\end{proof}

%===================================================================
\section{L'équation de Cauchy et la caractérisation de l'exponentielle}
%===================================================================

Le lemme suivant est le c\oe ur analytique de tout le reste. On ne suppose ni
continuité ni densité : la seule monotonie suffit.

\begin{lemme}[Cauchy multiplicatif, version monotone]\label{lem:cauchy}
Soit $\barF:[0,\infty)\to[0,1]$ décroissante, continue à droite en $0$, telle que
$\barF(0)=1$ et
\[
\barF(s+t)=\barF(s)\,\barF(t)\qquad\text{pour tous } s,t\ge0 .
\tag{$\star$}
\]
Si $\barF$ n'est pas identiquement égale à $1$ sur $[0,\infty)$, alors il existe
un unique $\lambda\in(0,\infty)$ tel que $\barF(t)=e^{-\lambda t}$ pour tout $t\ge0$.
\end{lemme}

\begin{proof}
\emph{Étape 1 : $\barF>0$ partout.} Supposons $\barF(t_0)=0$ pour un $t_0>0$.
Par ($\star$), $\barF(t_0)=\barF(t_0/2)^2$, donc $\barF(t_0/2)=0$, et par
récurrence $\barF(t_0/2^m)=0$ pour tout $m\in\N$. En faisant $m\to\infty$, la
continuité à droite en $0$ impose $1=\barF(0)=\lim_{m}\barF(t_0/2^m)=0$,
absurde. Donc $\barF(t)>0$ pour tout $t\ge0$.

\emph{Étape 2 : additivité de $\varphi=-\log\barF$.} L'étape 1 rend
$\varphi(t)=-\log\barF(t)\in[0,\infty)$ bien définie et finie. Par ($\star$),
$\varphi(s+t)=\varphi(s)+\varphi(t)$ ; de plus $\varphi$ est croissante (car
$\barF$ décroît) et $\varphi(0)=0$.

\emph{Étape 3 : $\varphi$ est linéaire sur $\Q_+$.} Posons $\lambda=\varphi(1)$.
Une récurrence immédiate sur l'additivité donne $\varphi(m t)=m\varphi(t)$ pour
tous $m\in\N$, $t\ge0$. Pour $p,q\in\N^*$ : d'une part
$\varphi(q\cdot\tfrac1q)=q\,\varphi(\tfrac1q)$, c'est-à-dire
$\varphi(\tfrac1q)=\tfrac\lambda q$ ; d'autre part
$\varphi(\tfrac pq)=p\,\varphi(\tfrac1q)=\tfrac pq\lambda$. Donc
$\varphi(r)=\lambda r$ pour tout $r\in\Q_+$.

\emph{Étape 4 : passage aux réels par encadrement.} Soit $t\ge0$ et $(r_m),(s_m)$
deux suites de rationnels positifs avec $r_m\le t\le s_m$, $r_m\uparrow t$ et
$s_m\downarrow t$. La croissance de $\varphi$ donne
$\lambda r_m=\varphi(r_m)\le\varphi(t)\le\varphi(s_m)=\lambda s_m$. En passant à
la limite, $\varphi(t)=\lambda t$, soit $\barF(t)=e^{-\lambda t}$.

\emph{Étape 5 : $\lambda\in(0,\infty)$ et unicité.} $\lambda=\varphi(1)\ge0$. Si
$\lambda=0$ alors $\barF\equiv1$, exclu par hypothèse ; $\lambda<\infty$ par
l'étape 1. L'unicité est claire puisque $\lambda=-\log\barF(1)$.
\end{proof}

\begin{theoreme}[Caractérisation de la loi exponentielle]\label{th:carac}
Soit $X>0$ p.s.\ de survie $\barF$ décroissante continue à droite,
$\barF(0)=1$, $X$ non p.s.\ infinie. Les assertions suivantes sont équivalentes.
\begin{enumerate}[label=(\roman*),leftmargin=2.2em,itemsep=0.15em]
  \item \textbf{(Absence de mémoire)} $\Prob(X>s+t\mid X>s)=\Prob(X>t)$ pour tous
  $s,t\ge0$ tels que $\barF(s)>0$.
  \item \textbf{(Survie multiplicative)} $\barF(s+t)=\barF(s)\barF(t)$ pour tous $s,t\ge0$.
  \item \textbf{(Loi)} Il existe $\lambda>0$ tel que $\barF(t)=e^{-\lambda t}$,
  c'est-à-dire $X\sim\Expo(\lambda)$.
  \item \textbf{(Taux constant)} $X$ admet une densité et son taux de panne
  vérifie $h\equiv\lambda$ sur $[0,\infty)$, pour un $\lambda>0$.
\end{enumerate}
\end{theoreme}

\begin{proof}
\emph{(i) $\Rightarrow$ (ii).} Montrons d'abord que (i) force $\barF>0$ partout.
Si $\barF(s_0)=0$ pour un $s_0$, soit $s^\ast=\inf\{s:\barF(s)=0\}$. Comme $X>0$
p.s.\ et $\barF$ est continue à droite en $0$ avec $\barF(0)=1$, on a
$s^\ast>0$. Pour $s<s^\ast$ on a $\barF(s)>0$ et (i) s'écrit
$\barF(s+t)=\barF(s)\barF(t)$ ; en prenant $s<s^\ast$ et $t<s^\ast$ avec
$s+t>s^\ast$ on obtient $0=\barF(s+t)=\barF(s)\barF(t)>0$, absurde (un tel
couple existe : prendre $s=t$ légèrement supérieur à $s^\ast/2$). Donc
$\barF>0$ partout, la restriction dans (i) est vide, et multiplier (i) par
$\barF(s)$ donne exactement (ii).

\emph{(ii) $\Rightarrow$ (iii).} C'est le lemme \ref{lem:cauchy} : $\barF$ n'est
pas identiquement $1$, sans quoi $X=+\infty$ p.s., ce qui est exclu.

\emph{(iii) $\Rightarrow$ (iv).} Si $\barF(t)=e^{-\lambda t}$, alors
$f(t)=\lambda e^{-\lambda t}$ est bien une densité, et
$h=f/\barF=\lambda e^{-\lambda t}/e^{-\lambda t}=\lambda$.

\emph{(iv) $\Rightarrow$ (iii).} Si $h\equiv\lambda$ alors $H(x)=\lambda x$ et le
théorème \ref{th:reconstruction} donne $\barF(x)=e^{-\lambda x}$.

\emph{(iii) $\Rightarrow$ (i).} Immédiat :
$\barF(s+t)/\barF(s)=e^{-\lambda(s+t)}/e^{-\lambda s}=e^{-\lambda t}=\barF(t)$.
\end{proof}

\begin{remarque}
Les quatre énoncés ne sont pas quatre faits liés mais un seul fait lu de quatre
façons : (ii) est (i) débarrassé du quotient, et (iv) est (ii) après passage au
logarithme, l'additivité de $H=-\log\barF$ équivalant à la constance de $H'=h$.
\end{remarque}

\begin{proposition}[Compensateur : le pont vers le cas général]\label{prop:compensateur}
Soit $X>0$ de hasard cumulé $H$ continu, strictement croissant sur $[0,\infty)$,
avec $H(0)=0$ et $H(x)\to+\infty$. Alors $H(X)\sim\Expo(1)$.
\end{proposition}

\begin{proof}
Sous ces hypothèses $H:[0,\infty)\to[0,\infty)$ est une bijection croissante
continue, d'inverse $H^{-1}$ croissante. Pour $t\ge0$,
\[
\Prob\bigl(H(X)>t\bigr)=\Prob\bigl(X>H^{-1}(t)\bigr)=\barF\bigl(H^{-1}(t)\bigr)
=e^{-H(H^{-1}(t))}=e^{-t},
\]
en utilisant le théorème \ref{th:reconstruction}. C'est la survie de $\Expo(1)$.
\end{proof}

Autrement dit : toute durée de vie devient sans mémoire une fois le temps
reparamétré par son propre hasard cumulé. L'exponentielle est le cas où cette
reparamétrisation est l'identité (à l'échelle $\lambda$ près).

\begin{center}\rule{0.35\textwidth}{0.4pt}\end{center}
\part*{Partie II --- Statistiques d'ordre}
\addcontentsline{toc}{section}{\textbf{Partie II --- Statistiques d'ordre}}

%===================================================================
\section{Cadre et densité jointe du vecteur ordonné}
%===================================================================

Soient $X_1,\dots,X_n$ i.i.d.\ de densité $f$ sur $(0,\infty)$. On note
$\Xord{1}\le\cdots\le\Xord{n}$ le réarrangement croissant et
$D_i=\Xord{i}-\Xord{i-1}$ les \emph{espacements}, avec $\Xord0=0$. On pose
\[
\Delta_n=\{x\in\R^n:\ 0<x_1<x_2<\cdots<x_n\}\quad\text{(ouvert de }\R^n).
\]

\begin{lemme}[Absence d'ex \ae quo]\label{lem:ties}
Si les $X_i$ sont i.i.d.\ de loi diffuse (sans atome), alors
$\Prob(\exists\, i\ne j:\ X_i=X_j)=0$. En particulier le vecteur ordonné est
p.s.\ à coordonnées deux à deux distinctes, et $(\Xord1,\dots,\Xord n)\in\Delta_n$ p.s.
\end{lemme}

\begin{proof}
Pour $i\ne j$, l'indépendance et Tonelli donnent
\[
\Prob(X_i=X_j)=\int_\R\Prob(X_i=y)\,\Prob_{X_j}(\dd y)=0,
\]
chaque $\Prob(X_i=y)$ étant nul par diffusion. On conclut par sous-additivité
sur les $\binom n2$ paires.
\end{proof}

\begin{theoreme}[Densité jointe des statistiques d'ordre]\label{th:densite-ordre}
Le vecteur $(\Xord1,\dots,\Xord n)$ admet la densité
\[
f_{(\cdot)}(x_1,\dots,x_n)=n!\,\prod_{i=1}^n f(x_i)\ \mathbf 1_{\Delta_n}(x_1,\dots,x_n).
\]
\end{theoreme}

\begin{proof}
Notons $\mathbf X_{(\cdot)}=(\Xord1,\dots,\Xord n)$ et, pour $\sigma\in\mathfrak S_n$,
\[
A_\sigma=\{X_{\sigma(1)}<X_{\sigma(2)}<\cdots<X_{\sigma(n)}\}.
\]
Par le lemme \ref{lem:ties}, les $A_\sigma$ sont deux à deux disjoints et
$\Prob\bigl(\bigcup_\sigma A_\sigma\bigr)=1$. Soit $B\subseteq\Delta_n$ borélien.
Sur $A_\sigma$, on a par définition
$\mathbf X_{(\cdot)}=(X_{\sigma(1)},\dots,X_{\sigma(n)})$, et de plus
$A_\sigma=\{(X_{\sigma(1)},\dots,X_{\sigma(n)})\in\Delta_n\}$. Comme
$B\subseteq\Delta_n$, l'inclusion $\{(X_{\sigma(1)},\dots,X_{\sigma(n)})\in B\}\subseteq A_\sigma$
est automatique, d'où
\[
\Prob\bigl(\mathbf X_{(\cdot)}\in B,\ A_\sigma\bigr)
=\Prob\bigl((X_{\sigma(1)},\dots,X_{\sigma(n)})\in B\bigr).
\]
Les $X_i$ étant i.i.d., elles sont échangeables : le vecteur
$(X_{\sigma(1)},\dots,X_{\sigma(n)})$ a la même loi que $(X_1,\dots,X_n)$. Donc
chacun des $n!$ termes vaut $\int_B\prod_{i}f(x_i)\dd x$. En sommant,
\[
\Prob\bigl(\mathbf X_{(\cdot)}\in B\bigr)
=\sum_{\sigma\in\mathfrak S_n}\Prob\bigl(\mathbf X_{(\cdot)}\in B,A_\sigma\bigr)
=n!\int_B\prod_{i=1}^n f(x_i)\dd x .
\]
Comme $\mathbf X_{(\cdot)}\in\Delta_n$ p.s., ceci caractérise la loi et identifie
la densité annoncée.
\end{proof}

%===================================================================
\section{L'application des espacements est un difféomorphisme}
%===================================================================

C'est l'ingrédient technique qui légitime le changement de variables. Soit
\[
\Phi:\R^n\to\R^n,\qquad
\Phi(d_1,\dots,d_n)=\Bigl(d_1,\ d_1+d_2,\ \dots,\ \sum_{k=1}^n d_k\Bigr),
\qquad\text{i.e.}\quad \Phi(d)_i=\sum_{k=1}^i d_k .
\]

\begin{lemme}[Difféomorphisme et jacobien]\label{lem:diffeo}
L'application $\Phi$ est linéaire, de matrice
$L=\bigl(\mathbf 1_{\{k\le i\}}\bigr)_{i,k}$, triangulaire inférieure à
diagonale unité. Elle vérifie :
\begin{enumerate}[label=(\alph*),leftmargin=2.2em,itemsep=0.15em]
  \item $\det L=1$, donc $\Phi$ est un automorphisme linéaire de $\R^n$ ;
  \item $\Phi$ réalise une bijection de $(0,\infty)^n$ sur $\Delta_n$, d'inverse
  $\Psi(x)=(x_1,\ x_2-x_1,\ \dots,\ x_n-x_{n-1})$ ;
  \item $\Phi|_{(0,\infty)^n}:(0,\infty)^n\to\Delta_n$ est un
  $C^\infty$-difféomorphisme entre ouverts de $\R^n$, de jacobien constant
  $\bigl|\det \Jac_\Phi(d)\bigr|=1$.
\end{enumerate}
\end{lemme}

\begin{proof}
\emph{(a)} $\Phi(d)_i=\sum_k L_{ik}d_k$ avec $L_{ik}=1$ si $k\le i$ et $0$
sinon : $\Phi$ est linéaire de matrice $L$, triangulaire inférieure de
coefficients diagonaux $L_{ii}=1$. Le déterminant d'une matrice triangulaire
étant le produit des coefficients diagonaux, $\det L=1\ne0$ et $L\in GL_n(\R)$.

\emph{(b)} Notons $\Psi$ l'application linéaire donnée par $\Psi(x)_1=x_1$ et
$\Psi(x)_k=x_k-x_{k-1}$ pour $k\ge2$. Vérifions $\Psi\circ\Phi=\mathrm{id}$ :
pour $k\ge2$,
$\Psi(\Phi(d))_k=\Phi(d)_k-\Phi(d)_{k-1}=\sum_{j\le k}d_j-\sum_{j\le k-1}d_j=d_k$,
et $\Psi(\Phi(d))_1=\Phi(d)_1=d_1$. De même $\Phi(\Psi(x))_i=\sum_{k\le i}\Psi(x)_k$
est une somme télescopique égale à $x_i$. Donc $\Psi=\Phi^{-1}$ sur $\R^n$.

Montrons maintenant les deux inclusions. Si $d\in(0,\infty)^n$, alors
$\Phi(d)_1=d_1>0$ et $\Phi(d)_i-\Phi(d)_{i-1}=d_i>0$, donc les coordonnées de
$\Phi(d)$ sont strictement croissantes et positives : $\Phi(d)\in\Delta_n$.
Réciproquement si $x\in\Delta_n$, alors $\Psi(x)_1=x_1>0$ et
$\Psi(x)_k=x_k-x_{k-1}>0$, donc $\Psi(x)\in(0,\infty)^n$. Ainsi
$\Phi\bigl((0,\infty)^n\bigr)\subseteq\Delta_n$ et
$\Psi(\Delta_n)\subseteq(0,\infty)^n$ ; comme $\Psi=\Phi^{-1}$, ces deux
inclusions donnent la bijectivité de $\Phi$ de $(0,\infty)^n$ sur $\Delta_n$.

\emph{(c)} $(0,\infty)^n$ et $\Delta_n$ sont ouverts dans $\R^n$ ($\Delta_n$ est
l'intersection finie des demi-espaces ouverts $\{x_1>0\}$ et
$\{x_i-x_{i-1}>0\}$). Une application linéaire est $C^\infty$, et $\Phi$ comme
$\Psi=\Phi^{-1}$ sont linéaires : la bijection de (b) est donc $C^\infty$ ainsi
que sa réciproque, c'est-à-dire un $C^\infty$-difféomorphisme. Enfin la
différentielle d'une application linéaire est elle-même en tout point, donc
$\Jac_\Phi(d)=L$ pour tout $d$ et $|\det \Jac_\Phi(d)|=|\det L|=1$.
\end{proof}

Nous utiliserons la forme suivante, standard, du transport de densité.

\begin{lemme}[Changement de variables pour les densités]\label{lem:cdv}
Soient $U,V\subseteq\R^n$ ouverts et $\Phi:U\to V$ un $C^1$-difféomorphisme. Si un
vecteur aléatoire $Z$ à valeurs dans $V$ admet la densité $g$, alors
$W=\Phi^{-1}(Z)$ admet sur $U$ la densité
$\ w\mapsto g\bigl(\Phi(w)\bigr)\,\bigl|\det \Jac_\Phi(w)\bigr|$.
\end{lemme}

\begin{proof}
Soit $B\subseteq U$ borélien. Alors $\{W\in B\}=\{Z\in\Phi(B)\}$ et $\Phi(B)$ est
borélien (image par un homéomorphisme). La formule du changement de variables
pour l'intégrale de Lebesgue, appliquée au difféomorphisme $\Phi$, donne
\[
\Prob(W\in B)=\int_{\Phi(B)}g(z)\dd z=\int_{B}g\bigl(\Phi(w)\bigr)\bigl|\det \Jac_\Phi(w)\bigr|\dd w .
\]
Ceci valant pour tout borélien $B\subseteq U$, l'intégrande est une densité de $W$.
\end{proof}

\begin{lemme}[Factorisation $\Rightarrow$ indépendance]\label{lem:factorisation}
Soit $W=(W_1,\dots,W_n)$ de densité $g$ sur $\R^n$. Si
$g(w)=\prod_{k=1}^n g_k(w_k)$ p.p., où chaque $g_k\ge0$ vérifie
$\int_\R g_k=1$, alors les $W_k$ sont indépendantes et $W_k$ a pour densité $g_k$.
\end{lemme}

\begin{proof}
Pour des boréliens $B_1,\dots,B_n$, Tonelli donne
\[
\Prob(W_1\in B_1,\dots,W_n\in B_n)=\int_{B_1\times\cdots\times B_n}\prod_k g_k(w_k)\dd w
=\prod_{k=1}^n\int_{B_k}g_k(w_k)\dd w_k .
\]
En prenant $B_j=\R$ pour $j\ne k$ et en utilisant $\int g_j=1$, on obtient
$\Prob(W_k\in B_k)=\int_{B_k}g_k$ : $g_k$ est la densité de $W_k$. L'égalité
ci-dessus devient alors
$\Prob(\bigcap_k\{W_k\in B_k\})=\prod_k\Prob(W_k\in B_k)$, qui est l'indépendance.
\end{proof}

%===================================================================
\section{Le théorème de Rényi}
%===================================================================

\begin{theoreme}[Rényi, 1953]\label{th:renyi}
Soient $X_1,\dots,X_n$ i.i.d.\ de loi $\Expo(\lambda)$, $\lambda>0$. Les
espacements $D_1,\dots,D_n$, où $D_i=\Xord i-\Xord{i-1}$ et $\Xord0=0$, sont
\emph{indépendants} et
\[
D_i\sim\Expo\bigl((n-i+1)\lambda\bigr),\qquad i=1,\dots,n .
\]
De façon équivalente, si $E_1,\dots,E_n$ sont i.i.d.\ de loi $\Expo(1)$, alors
\[
\bigl(\Xord1,\dots,\Xord n\bigr)\ \eqd\
\Bigl(\frac1\lambda\sum_{k=1}^{i}\frac{E_k}{\,n-k+1\,}\Bigr)_{1\le i\le n}
\qquad\text{(égalité en loi des vecteurs).}
\]
\end{theoreme}

\begin{proof}
\emph{Étape 1 : densité du vecteur ordonné.} Par le théorème
\ref{th:densite-ordre} avec $f(x)=\lambda e^{-\lambda x}\mathbf 1_{\{x>0\}}$, le
vecteur $\mathbf X_{(\cdot)}$ a pour densité sur $\Delta_n$
\[
f_{(\cdot)}(x)=n!\,\lambda^n\exp\Bigl(-\lambda\sum_{i=1}^n x_i\Bigr),\qquad x\in\Delta_n .
\]

\emph{Étape 2 : transport.} Le vecteur des espacements est exactement
$D=\Psi(\mathbf X_{(\cdot)})=\Phi^{-1}(\mathbf X_{(\cdot)})$, avec $\Phi,\Psi$ du
lemme \ref{lem:diffeo}. Ce lemme fournit un $C^\infty$-difféomorphisme
$\Phi:(0,\infty)^n\to\Delta_n$ de jacobien $|\det \Jac_\Phi|\equiv1$. Le lemme
\ref{lem:cdv} donne donc la densité de $D$ sur $(0,\infty)^n$ :
\[
g(d)=f_{(\cdot)}\bigl(\Phi(d)\bigr)\cdot\underbrace{\bigl|\det \Jac_\Phi(d)\bigr|}_{=1}
=n!\,\lambda^n\exp\Bigl(-\lambda\sum_{i=1}^n \Phi(d)_i\Bigr).
\]

\emph{Étape 3 : l'identité de sommation triangulaire.} Il s'agit d'un simple
échange de sommations finies (Fubini discret) :
\[
\sum_{i=1}^n\Phi(d)_i=\sum_{i=1}^n\sum_{k=1}^i d_k
=\sum_{i=1}^n\sum_{k=1}^n d_k\,\mathbf 1_{\{k\le i\}}
=\sum_{k=1}^n d_k\sum_{i=1}^n\mathbf 1_{\{i\ge k\}}
=\sum_{k=1}^n (n-k+1)\,d_k ,
\]
puisque $\#\{i:\ k\le i\le n\}=n-k+1$.

\emph{Étape 4 : factorisation.} En posant $\lambda_k=(n-k+1)\lambda$ et en
utilisant $\prod_{k=1}^n(n-k+1)=n!$ :
\[
g(d)=n!\,\lambda^n\exp\Bigl(-\sum_{k=1}^n\lambda_k d_k\Bigr)
=\prod_{k=1}^n\lambda_k\cdot\prod_{k=1}^n e^{-\lambda_k d_k}
=\prod_{k=1}^n\Bigl[\lambda_k e^{-\lambda_k d_k}\Bigr],\qquad d\in(0,\infty)^n .
\]
Chaque facteur $g_k(u)=\lambda_k e^{-\lambda_k u}\mathbf 1_{\{u>0\}}$ est une
densité de probabilité ($\int_0^\infty\lambda_k e^{-\lambda_k u}\dd u=1$). Le
lemme \ref{lem:factorisation} conclut : les $D_k$ sont indépendantes de lois
respectives $\Expo(\lambda_k)=\Expo((n-k+1)\lambda)$.

\emph{Étape 5 : forme de Rényi.} Posons $E_k=\lambda_k D_k$. Pour $t\ge0$,
$\Prob(E_k>t)=\Prob(D_k>t/\lambda_k)=e^{-\lambda_k\cdot t/\lambda_k}=e^{-t}$,
donc $E_k\sim\Expo(1)$, et les $E_k$ restent indépendantes comme images
mesurables de variables indépendantes. Enfin
\[
\Xord i=\sum_{k=1}^i D_k=\sum_{k=1}^i\frac{E_k}{\lambda_k}
=\frac1\lambda\sum_{k=1}^i\frac{E_k}{n-k+1},
\]
cette fois comme identité \emph{presque sûre} sur l'espace de départ, ce qui
entraîne l'égalité en loi des vecteurs annoncée. \qedhere
\end{proof}

\begin{remarque}
Le point où l'exponentialité intervient est l'étape 2--4, et uniquement là : la
densité jointe ne dépend de $x$ que par $\sum_i x_i$, forme linéaire que le
changement de variables triangulaire transforme en une combinaison linéaire à
coefficients \emph{séparés} $\sum_k(n-k+1)d_k$. C'est la traduction analytique de
l'absence de mémoire (théorème \ref{th:carac}(ii)) : seule une survie
multiplicative fait apparaître une exponentielle d'une somme, donc un produit.
Pour une loi générale la densité jointe reste $n!\prod f(x_i)$, mais
$\prod_i f(\Phi(d)_i)$ ne se factorise pas selon les $d_k$ et les espacements
sont en général dépendants.
\end{remarque}

%===================================================================
\section{La preuve probabiliste : compétition et renouvellement}
%===================================================================

La deuxième démonstration explique \emph{pourquoi} le résultat est vrai. Sa
difficulté est qu'elle semble conditionner par le temps aléatoire $\Xord1$. On
l'évite entièrement par le lemme suivant, dont le calcul livre d'un coup toutes
les indépendances.

\begin{lemme}[Compétition d'exponentielles et résidus]\label{lem:competition}
Soient $Y_1,\dots,Y_m$ indépendantes, $Y_i\sim\Expo(\mu_i)$, et
$\Lambda=\sum_{i=1}^m\mu_i$. Posons
\[
M=\min_{1\le i\le m}Y_i,\qquad I=\argmin_{1\le i\le m}Y_i\quad(\text{p.s.\ bien défini}),
\qquad R_i=Y_i-M\ \ (i\ne I).
\]
Alors, pour tout $j$, tous $t\ge0$ et tous $r_i\ge0$ $(i\ne j)$ :
\[
\Prob\Bigl(I=j,\ M>t,\ \bigcap_{i\ne j}\{R_i>r_i\}\Bigr)
=\frac{\mu_j}{\Lambda}\;\cdot\;e^{-\Lambda t}\;\cdot\;\prod_{i\ne j}e^{-\mu_i r_i}.
\tag{$\dagger$}
\]
Par conséquent :
$\Prob(I=j)=\mu_j/\Lambda$ ; $M\sim\Expo(\Lambda)$ ; conditionnellement à
$\{I=j\}$, les $m$ variables $M$ et $(R_i)_{i\ne j}$ sont indépendantes, avec
$R_i\sim\Expo(\mu_i)$ ; et le couple $(I,M)$ est indépendant de la famille des
résidus.
\end{lemme}

\begin{proof}
Le lemme \ref{lem:ties} garantit l'unicité p.s.\ du minimiseur, donc $I$ est
p.s.\ bien défini. Fixons $j$, $t\ge0$ et $r_i\ge0$. Observons l'identité
d'événements, à un ensemble négligeable près :
\[
\Bigl\{I=j,\ M>t,\ \bigcap_{i\ne j}\{Y_i-M>r_i\}\Bigr\}
=\Bigl\{Y_j>t,\ \bigcap_{i\ne j}\{Y_i>Y_j+r_i\}\Bigr\}.
\]
En effet, sur le membre de droite on a $Y_i>Y_j+r_i\ge Y_j$ pour tout $i\ne j$,
donc $M=Y_j$ et $I=j$, puis $Y_i-M>r_i$ ; réciproquement, sur le membre de gauche
$M=Y_j$ et les inégalités se réécrivent. Par indépendance, en intégrant selon la
loi de $Y_j$ (Tonelli, intégrande positive) :
\[
\Prob\Bigl(Y_j>t,\bigcap_{i\ne j}\{Y_i>Y_j+r_i\}\Bigr)
=\int_t^{\infty}\mu_j e^{-\mu_j y}\prod_{i\ne j}\Prob\bigl(Y_i>y+r_i\bigr)\dd y .
\]
Or $\Prob(Y_i>y+r_i)=e^{-\mu_i(y+r_i)}=e^{-\mu_i y}e^{-\mu_i r_i}$ --- c'est ici,
et seulement ici, qu'intervient l'absence de mémoire, sous sa forme
multiplicative $\barF(y+r)=\barF(y)\barF(r)$ du théorème \ref{th:carac}(ii). En
sortant les facteurs constants en $y$ :
\[
=\Bigl(\prod_{i\ne j}e^{-\mu_i r_i}\Bigr)\int_t^{\infty}\mu_j
\exp\Bigl(-y\bigl(\mu_j+\textstyle\sum_{i\ne j}\mu_i\bigr)\Bigr)\dd y
=\Bigl(\prod_{i\ne j}e^{-\mu_i r_i}\Bigr)\,\mu_j\int_t^\infty e^{-\Lambda y}\dd y ,
\]
et $\int_t^\infty e^{-\Lambda y}\dd y=e^{-\Lambda t}/\Lambda$, ce qui donne
exactement ($\dagger$).

Les conséquences se lisent sur ($\dagger$) en spécialisant. En prenant $t=0$ et
tous les $r_i=0$ : $\Prob(I=j)=\mu_j/\Lambda$ (et ces quantités somment bien à
$1$). En sommant ($\dagger$) sur $j$ avec $r_i\equiv0$ : $\Prob(M>t)=e^{-\Lambda t}$,
soit $M\sim\Expo(\Lambda)$. Enfin ($\dagger$) divisé par $\Prob(I=j)$ s'écrit
\[
\Prob\Bigl(M>t,\bigcap_{i\ne j}\{R_i>r_i\}\ \Bigm|\ I=j\Bigr)
=e^{-\Lambda t}\prod_{i\ne j}e^{-\mu_i r_i},
\]
produit de fonctions de survie : les queues jointes se factorisent sur le pavé
$[0,\infty)^{m}$, ce qui caractérise l'indépendance de $M$ et des $(R_i)_{i\ne j}$
sous $\Prob(\cdot\mid I=j)$, avec les marginales annoncées. Comme les lois
conditionnelles des résidus ne dépendent pas de $j$ ni de $t$, la famille
$(R_i)$ est indépendante de $(I,M)$.
\end{proof}

\begin{proof}[Seconde démonstration du théorème \ref{th:renyi}]
Récurrence sur $n$. Pour $n=1$, $D_1=X_{(1)}=X_1\sim\Expo(\lambda)$ et l'énoncé
est vrai.

Supposons le résultat acquis au rang $n-1$. Appliquons le lemme
\ref{lem:competition} à $m=n$ et $\mu_i=\lambda$, de sorte que $\Lambda=n\lambda$.
Il donne : $D_1=\Xord1=M\sim\Expo(n\lambda)$ ; et, conditionnellement à
$\{I=j\}$, les résidus $(R_i)_{i\ne j}$ sont i.i.d.\ $\Expo(\lambda)$ et
indépendants de $M$. Comme cette loi conditionnelle ne dépend pas de $j$, on
conclut inconditionnellement : la famille $R=(R_i)_{i\ne I}$ est composée de
$n-1$ variables i.i.d.\ $\Expo(\lambda)$, indépendantes de $D_1$.

Or les statistiques d'ordre des $R_i$ sont exactement les
$\Xord{k+1}-\Xord1$ pour $k=1,\dots,n-1$ : le réarrangement croissant commute
avec la translation par la constante $M$ et la suppression du minimum.
Autrement dit, les espacements de l'échantillon $R$ sont $D_2,\dots,D_n$.
L'hypothèse de récurrence appliquée à $R$ (échantillon i.i.d.\ $\Expo(\lambda)$
de taille $n-1$) fournit alors : $D_2,\dots,D_n$ indépendantes avec
\[
D_{i}\sim\Expo\bigl(((n-1)-(i-1)+1)\lambda\bigr)=\Expo\bigl((n-i+1)\lambda\bigr),
\qquad i=2,\dots,n .
\]
Enfin $(D_2,\dots,D_n)$ est fonction mesurable de $R$, donc indépendante de
$D_1$. Les $n$ espacements sont donc globalement indépendants avec les lois
voulues.
\end{proof}

\begin{remarque}
La rigueur tient au fait que ($\dagger$) est une égalité de lois \emph{jointes}
calculée d'avance : à aucun moment on ne conditionne par la valeur du temps
aléatoire $\Xord1$. L'énoncé intuitif « à l'instant de la première panne, les
$n-1$ horloges survivantes redémarrent à neuf » est précisément le contenu de la
factorisation du membre de droite de ($\dagger$) en
$\Prob(I=j)\times\Prob(M>t)\times\prod_i\Prob(R_i>r_i)$.
\end{remarque}

%===================================================================
\section{Conséquences}
%===================================================================

On note $H_m=\sum_{j=1}^m 1/j$ le $m$-ième nombre harmonique, $H_0=0$.

\begin{corollaire}[Moments]
Pour $1\le i\le n$,
\[
\E[\Xord i]=\frac{1}{\lambda}\bigl(H_n-H_{n-i}\bigr),
\qquad
\Var(\Xord i)=\frac{1}{\lambda^2}\sum_{j=n-i+1}^{n}\frac{1}{j^2}.
\]
En particulier $\E[\Xord n]=H_n/\lambda$ et
$\Var(\Xord n)=\lambda^{-2}\sum_{j=1}^n j^{-2}\to\pi^2/(6\lambda^2)$.
\end{corollaire}

\begin{proof}
Par le théorème \ref{th:renyi}, $\Xord i=\sum_{k=1}^i D_k$ avec $D_k$
indépendantes, $\E[D_k]=1/((n-k+1)\lambda)$ et
$\Var(D_k)=1/((n-k+1)\lambda)^2$. La linéarité de l'espérance et l'additivité de
la variance \emph{sous indépendance} donnent
\[
\E[\Xord i]=\frac1\lambda\sum_{k=1}^i\frac{1}{n-k+1},
\qquad
\Var(\Xord i)=\frac1{\lambda^2}\sum_{k=1}^i\frac{1}{(n-k+1)^2}.
\]
Le changement d'indice $j=n-k+1$ (bijection de $\{1,\dots,i\}$ sur
$\{n-i+1,\dots,n\}$) transforme la première somme en
$\sum_{j=n-i+1}^n 1/j=H_n-H_{n-i}$ et la seconde en $\sum_{j=n-i+1}^n 1/j^2$.
Pour $i=n$ on retrouve $H_n$ et $\sum_{j=1}^n 1/j^2$, dont la limite est
$\zeta(2)=\pi^2/6$.
\end{proof}

\begin{theoreme}[Domaine d'attraction : loi de Gumbel]
Avec $X_i$ i.i.d.\ $\Expo(\lambda)$,
\[
\lambda\Xord n-\log n\ \cvloi\ G,\qquad \Prob(G\le x)=e^{-e^{-x}},\ x\in\R .
\]
\end{theoreme}

\begin{proof}
Fixons $x\in\R$ et soit $n>e^{-x}$, de sorte que $(x+\log n)/\lambda>0$. Comme
$\Xord n=\max_i X_i$, l'indépendance donne
\[
\Prob\bigl(\lambda\Xord n-\log n\le x\bigr)
=\Prob\Bigl(\max_i X_i\le\frac{x+\log n}{\lambda}\Bigr)
=\Bigl[F\Bigl(\frac{x+\log n}{\lambda}\Bigr)\Bigr]^{n}
=\Bigl[1-e^{-(x+\log n)}\Bigr]^{n},
\]
soit $\bigl(1-\frac{e^{-x}}{n}\bigr)^{n}$. Or pour $a\in\R$ fixé,
$n\log(1-a/n)=n\bigl(-a/n+O(n^{-2})\bigr)\to-a$, donc
$(1-a/n)^n\to e^{-a}$. Avec $a=e^{-x}$ :
\[
\Prob\bigl(\lambda\Xord n-\log n\le x\bigr)\xrightarrow[n\to\infty]{}e^{-e^{-x}} .
\]
La limite est une fonction de répartition continue sur $\R$, donc la convergence
en tout point équivaut à la convergence en loi.
\end{proof}

\begin{remarque}
La démonstration n'utilise que $\barF(x)=e^{-\lambda x}$ exactement, c'est-à-dire
la queue exponentielle pure fournie par le taux de panne constant. Elle
n'utilise \emph{pas} le théorème \ref{th:renyi} : l'argument souvent avancé
(« $\E[\Xord n]\sim\log n/\lambda$ et $\Var(\Xord n)$ reste bornée, donc la
limite existe ») ne constitue pas une preuve, la tension d'une suite ne
suffisant pas à identifier une loi limite.
\end{remarque}

\vspace{1.2em}
\begin{center}
\rule{0.35\textwidth}{0.4pt}\\[0.5em]
\textit{Le taux de panne est l'opposé de la dérivée logarithmique de la survie ;
l'exponentielle est l'unique loi de taux constant ; et cette constance, lue au
niveau du vecteur ordonné, se traduit exactement par l'indépendance des
espacements.}
\end{center}

\end{document}
